\documentclass[11pt,a4paper]{article}
\usepackage[T1]{fontenc}
\usepackage{lmodern}
\usepackage[margin=25mm]{geometry}
\usepackage{amsmath,amssymb,amsthm}
\usepackage{microtype}
\usepackage[colorlinks=true,linkcolor=blue,citecolor=blue,urlcolor=blue]{hyperref}
\hypersetup{pdftitle={Inverse Frobenius in generic multiplicative fields},
 pdfauthor={Alec Kriebel},pdflang={en},
 pdfsubject={A negative answer to the unrestricted integral endomorphism question}}
\newtheorem{theorem}{Theorem}
\newtheorem{lemma}[theorem]{Lemma}
\newtheorem{corollary}[theorem]{Corollary}
\newcommand{\Z}{\mathbb Z}
\newcommand{\End}{\operatorname{End}_{\mathrm{def}}}
\setlength{\emergencystretch}{1em}
\title{Inverse Frobenius in generic multiplicative fields}
\author{Alec Kriebel\\Independent researcher\\
 \small\href{https://orcid.org/0009-0001-9320-500X}{ORCID: 0009-0001-9320-500X}}
\date{October 13, 2026 \quad Version 1.0}
\begin{document}
\maketitle
\begin{abstract}
In every positive-characteristic model of ACFH, inverse Frobenius is a
parameter-free definable multiplicative endomorphism outside the integral
polynomial ring generated by the generic operator. This gives a negative
answer to the unrestricted equality stated in the 2023 Oberwolfach
contribution and in arXiv:2212.02115v4, Question 5.8.
We include an elementary existential-closure proof of polynomial
faithfulness. The characteristic-zero case and the full classification
after adjoining inverse Frobenius remain unresolved here.
\end{abstract}

\section{The question and the exact answer}
Let $T$ be the theory of algebraically closed fields $K$ with a unary
operator $\theta$, where $\theta(0)=0$ and
$\theta|_{K^\times}:K^\times\to K^\times$ is a group endomorphism.
The models of ACFH are precisely the existentially closed models of $T$
\cite[Corollary 3.15]{de25}. Such models exist in every prime
characteristic: apply the model-companion embedding
\cite[Corollary 3.16]{de25} to the algebraic closure of $\mathbb F_p$
with the identity operator. The embedding preserves characteristic.

Write $E=\End(K^\times;\theta)$ for the ring of group endomorphisms
definable in the field language expanded by $\theta$.
Its addition is pointwise multiplication and its multiplication is
composition. Its zero is the constant-one map; its identity is $x\mapsto x$.
The integer $n$ denotes $[n](x)=x^n$, with ordinary group inverses when
$n<0$. For $P(X)=\sum_{i=0}^d a_iX^i\in\Z[X]$, put
\[
 P(\theta)(x)=\prod_{i=0}^d\theta^i(x)^{a_i},\qquad x\in K^\times.
\]
The image of evaluation is denoted $\Z[\theta]$.
These coefficients remain integers even when the field has characteristic
$p$: the operator $[p]$ is Frobenius, rather than the ring zero.

The question whether $E=\Z[\theta]$ is stated without a characteristic
restriction in the Oberwolfach contribution \cite[pp.~97--98]{ow23}
and in \cite[Question 5.8]{de25}. We answer that exact unrestricted
question negatively.

\begin{theorem}\label{thm:main}
Let $(K,\theta)$ be a model of ACFH of characteristic $p>0$.
The map $\rho:K\to K$ defined by $\rho(x)^p=x$ is a parameter-free
definable field automorphism commuting with $\theta$.
Its restriction to $K^\times$ lies in $E\setminus\Z[\theta]$.
\end{theorem}

The inverse-Frobenius map is classical; negative Frobenius powers also
appear in \cite[Proposition 1.14]{de19}. Polynomial faithfulness is
already \cite[Corollary 5.7]{de25}. The counterexample below combines
these known ingredients. The next lemma supplies a direct derivation of
the latter ingredient, without assuming an endomorphism classification.

\section{A checkable proof}
We use the elementary extension property of divisible abelian groups:
if $H\leq A$ and $D$ is divisible, any homomorphism $H\to D$ extends
to $A$. Here is the needed proof. Take a maximal partial extension
$h:B\to D$ by Zorn's lemma. If $a\in A\setminus B$, then
$\{k\in\Z:a^k\in B\}$ is either $\{0\}$ or $n\Z$ for $n>0$.
In the first case choose any $d\in D$; in the second choose $d$ with
$d^n=h(a^n)$. The formula $h(ba^k)=h(b)d^k$ defines a homomorphism
on $B\langle a\rangle$: uniqueness of the expression handles the first
case, and the relation subgroup $n\Z$ handles the second.
This contradicts maximality, so $B=A$.

\begin{lemma}\label{lem:faithful}
If $(K,\theta)$ is existentially closed in $T$, evaluation
$\Z[X]\to E$ at $\theta$ is injective, in every characteristic.
\end{lemma}
\begin{proof}
Let $Q(X)=\sum_{i=0}^d b_iX^i\ne0$.
Choose algebraically independent $t_0,\ldots,t_d$ over $K$ and set
$L=\overline{K(t_0,\ldots,t_d)}$.
The subgroup
\[
 H=K^\times\langle t_0,\ldots,t_{d-1}\rangle\leq L^\times
\]
is the direct product of $K^\times$ with the free abelian group on the
displayed indeterminates. A nontrivial Laurent-monomial relation over
$K$ would contradict algebraic independence. On $H$ define
\[
 h\left(c\prod_{i<d}t_i^{m_i}\right)
       =\theta(c)\prod_{i<d}t_{i+1}^{m_i}.
\]
For $d=0$, use $H=K^\times$ and $h=\theta|_{K^\times}$.
This is a well-defined homomorphism. Since $L^\times$ is divisible,
the extension property gives a homomorphism $\theta_L:L^\times\to
L^\times$ extending $h$. Set $\theta_L(0)=0$.
Thus $(L,\theta_L)$ is a $T$-extension, and
$\theta_L^i(t_0)=t_i$ for $0\leq i\leq d$. Consequently
\[
 Q(\theta_L)(t_0)=\prod_{i=0}^d t_i^{b_i}\ne1.
\]
The transferred existential formula can be written without negative
powers or an inverse symbol:
\[
 \exists x\left(x\ne0\ \wedge\
 \prod_{b_i>0}\theta^i(x)^{b_i}\ne
 \prod_{b_i<0}\theta^i(x)^{-b_i}\right).
\]
Empty products are $1$. It holds in the $T$-extension, hence in $K$
by existential closedness. Therefore $Q(\theta)$ is not the zero
endomorphism. Constants and negative coefficients are included.
\end{proof}

\begin{proof}[Proof of Theorem~\ref{thm:main}]
Algebraic closedness makes $x\mapsto x^p$ surjective.
It is injective because $u^p-v^p=(u-v)^p$.
Its inverse $\rho$ has the parameter-free graph $y^p=x$.
Unique roots give
\[
 \rho(xz)=\rho(x)\rho(z),\qquad
 \rho(x+z)=\rho(x)+\rho(z).
\]
Thus $\rho$ is a field automorphism. Multiplicativity of $\theta$ gives
$\theta(\rho(x))^p=\theta(x)$, so uniqueness also gives
$\theta\rho=\rho\theta$, including at zero.

Suppose that $\rho|_{K^\times}=P(\theta)$ for $P\in\Z[X]$.
Compose with $[p]$ to obtain
\[
 (pP-1)(\theta)=0_E.
\]
Lemma~\ref{lem:faithful} forces $pP-1=0$ in $\Z[X]$.
If $a_0$ is the constant coefficient of $P$, the constant coefficient
of $pP-1$ is $pa_0-1$, nonzero modulo $p$.
This is impossible.
\end{proof}

\begin{corollary}\label{cor:localized}
In characteristic $p>0$, the subring of $E$ generated by
$\Z[\theta]$ and $\rho$ is isomorphic to $\Z[1/p][X]$ and
strictly contains $\Z[\theta]$.
\end{corollary}
\begin{proof}
The scalar $[p]$ commutes with every group endomorphism, since
$f(x^p)=f(x)^p$. Its inverse $\rho$ is therefore central.
Define $P/p^n\mapsto\rho^n\circ P(\theta)$.
Clearing common powers of $p$, which are central units in $E$,
proves well-definedness and compatibility with the ring operations.
Composing an element of its kernel with $[p^n]$ gives
$P(\theta)=0_E$, hence $P=0$ by Lemma~\ref{lem:faithful}.
Strictness follows from Theorem~\ref{thm:main}.
\end{proof}

\section{Scope, prior work, and verification}
Theorem~\ref{thm:main} completely refutes the unrestricted equality;
Corollary~\ref{cor:localized} is only an inclusion.
Equality with $\Z[1/p][\theta]$ in characteristic $p$ and the
characteristic-zero classification are separate questions, unresolved
here. Existential closedness is essential: with $\theta=\rho$ on
an arbitrary algebraically closed characteristic-$p$ field,
$pX-1$ annihilates $\theta$ and $\rho\in\Z[\theta]$.
In finite fields inverse Frobenius is an integer power; finite-field
experiments therefore cannot establish the theorem's exclusion.

Chini's vector-space classification and subsequent question distinguish
the operator reduct from a richer base language that already defines
additional endomorphisms \cite[Theorem 4.10 and Question 4.11]{chini3}.
This is relevant prior conceptual context. The proposed adaptation to
multiplicative fields explicitly involves torsion
\cite[Remark 2.20]{chini1}. On $K^\times/\operatorname{Tor}(K^\times)$,
inverse Frobenius induces the existing rational scalar $1/p$;
the present exclusion concerns the full group and integral coefficients.
Those vector-space results do not supply that full-group exclusion.

A bounded primary-source and repository audit found no earlier exact
full counterexample; first historical priority is not certified.
The final APAL journal body was inaccessible and was not compared with
the inspected arXiv v4. Our negative answer targets the explicitly
identified OWR and arXiv-v4 formulations; no assertion about the final
journal wording or the question's current open status is made.
The supporting archive contains the frozen candidate, proof and source
audits, priority evidence and clarification, and portable integrity and
finite negative controls. Those computations check their stated finite
identities and package integrity, rather than formally verify the theorem.

\paragraph{Disclosure.}
AI tools were used extensively in developing, reproducing, checking,
adversarially reviewing, and preparing this result and its materials.
This is an unrefereed research note. Internal AI-assisted review is
neither external human peer review nor formal proof-assistant
certification. The deposit is licensed CC BY 4.0.

\begin{thebibliography}{9}
\small\raggedright
\bibitem{ow23} C.~d'Elb\'ee,
\emph{Generic multiplicative endomorphism of fields}, contribution in
\emph{Model Theory: Combinatorics, Groups, Valued Fields and Neostability},
Oberwolfach Reports \textbf{20} (2023), Report 2, pp.~97--98.
\href{https://doi.org/10.4171/OWR/2023/2}{doi:10.4171/OWR/2023/2}.

\bibitem{de25} C.~d'Elb\'ee,
\emph{Generic multiplicative endomorphism of a field},
\href{https://arxiv.org/abs/2212.02115v4}{arXiv:2212.02115v4}
(revision January 17, 2025), Sections 3.3 and 5.1.
Published as Annals of Pure and Applied Logic \textbf{176}(4) (2025),
103554, \href{https://doi.org/10.1016/j.apal.2025.103554}{doi:10.1016/j.apal.2025.103554}.
The version inspected for this note is arXiv v4.

\bibitem{de19} C.~d'Elb\'ee,
\emph{Forking, Imaginaries and other features of ACFG},
\href{https://arxiv.org/abs/1812.09378v2}{arXiv:1812.09378v2}
(revision October 30, 2019), Proposition 1.14.

\bibitem{chini1} L.~Chini,
\emph{Model Theory of Generic Vector Space Endomorphisms},
\href{https://arxiv.org/abs/2502.13667v3}{arXiv:2502.13667v3}
(revision July 15, 2026), Remark 2.20.

\bibitem{chini3} L.~Chini,
\emph{Model theory of generic vector space endomorphisms III: Reducts},
\href{https://arxiv.org/abs/2607.17014v1}{arXiv:2607.17014v1}
(submitted July 19, 2026), Section 4.2.
\end{thebibliography}
\end{document}
